Observe that
$\ang{1} = \ang{21} - \ang{20}$,
and $\ang{20} = \mfrac{\pi}{9}$.
Hence:
\begin{align*}
 \sin\paren{\ang{1}}
 &= \sin\paren{\ang{21} - \ang{20}}
  = \sin\paren{\ang{21}} \cos\paren{\ang{20}}
    - \cos\paren{\ang{21}} \sin\paren{\ang{20}} = \\
 &= \frac{1}{2} \sin\paren{\ang{21}}
    \paren{e^{i \frac{\pi}{9}} + e^{ -i \frac{\pi}{9}}}
    - \frac{1}{2 i} \cos\paren{\ang{21}}
    \paren{e^{i \frac{\pi}{9}} - e^{ -i \frac{\pi}{9}}} = \\
 &= \frac{1}{2} e^{i \frac{\pi}{9}}
    \paren{ \sin\paren{\ang{21}} + i \cos\paren{\ang{21}}}
    + \frac{1}{2} e^{ -i \frac{\pi}{9}}
    \paren{ \sin\paren{\ang{21}} - i \cos\paren{\ang{21}}}
\end{align*}
and:
\begin{align*}
 \cos\paren{\ang{1}}
 &= \cos\paren{\ang{21} - \ang{20}}
  = \cos\paren{\ang{21}} \cos\paren{\ang{20}}
    + \sin\paren{\ang{21}} \sin\paren{\ang{20}} = \\
 &= \frac{1}{2} \cos\paren{\ang{21}}
    \paren{e^{i \frac{\pi}{9}} + e^{ -i \frac{\pi}{9}}}
    + \frac{1}{2 i} \sin\paren{\ang{21}}
    \paren{e^{i \frac{\pi}{9}} - e^{ -i \frac{\pi}{9}}} = \\
 &= \frac{1}{2} e^{i \frac{\pi}{9}}
    \paren{ \cos\paren{\ang{21}} - i \sin\paren{\ang{21}}}
    + \frac{1}{2} e^{ -i \frac{\pi}{9}}
    \paren{ \cos\paren{\ang{21}} + i \sin\paren{\ang{21}}}
\end{align*}
The term $e^{i \frac{\pi}{9}}$
is the principal cube root of
$e^{i \frac{\pi}{3}} = \mfrac{1}{2} + i \mfrac{\sqrt{3}}{2}$,
and the term $e^{- i \frac{\pi}{9}}$
is the principal cube root of
$e^{- i \frac{\pi}{3}} = \mfrac{1}{2} - i \mfrac{\sqrt{3}}{2}$,
so they can be written as:
\begin{align*}
 e^{i \frac{\pi}{9}}
 = \sqrt[3]{\frac{1 + i \sqrt{3}}{2}} \,\, , \,\,\,\,\,\,
 e^{- i \frac{\pi}{9}}
 = \sqrt[3]{\frac{1 - i \sqrt{3}}{2}}
\end{align*}
The values of $\sin\paren{\ang{21}}$ and $\cos\paren{\ang{21}}$ are the following (they can be calculated using the angle difference formulas from $\ang{21} = \ang{36} - \ang{15}$):
\begin{align*}
 \sin\paren{\ang{21}}
 &= \frac{1}{8} \big( \sqrt{3} + 1 \big) \sqrt{5 - \sqrt{5}}
     - \frac{1}{16} \sqrt{2} \, \big( \sqrt{3} - 1 \big) \big(\sqrt{5} + 1 \big)
 \\
 \cos\paren{\ang{21}}
 &= \frac{1}{8} \big( \sqrt{3} - 1 \big) \sqrt{5 - \sqrt{5}}
     + \frac{1}{16} \sqrt{2} \, \big( \sqrt{3} + 1 \big) \big( \sqrt{5} + 1 \big)
\end{align*}
Altogether:
\begin{align*}
 \sin\paren{\ang{1}}
 = \frac{1}{16} \, \sqrt[3]{\frac{1 + i \sqrt{3}}{2}}
 \, \cdot \,
   &\Bigg( \!
     \big( \sqrt{3} + 1 \big) \sqrt{5 - \sqrt{5}}
     - \frac{1}{2} \sqrt{2} \, \big( \sqrt{3} - 1 \big) \big( \sqrt{5} + 1 \big) \, +
   \\
   &\,\,\,\,
     + \, i \, \bigg( \!
     \big( \sqrt{3} - 1 \big) \sqrt{5 - \sqrt{5}}
     + \frac{1}{2} \sqrt{2} \, \big( \sqrt{3} + 1 \big) \big( \sqrt{5} + 1 \big) \!
     \bigg) \!
   \Bigg) \,\, +
 \\
 + \, \frac{1}{16} \, \sqrt[3]{\frac{1 - i \sqrt{3}}{2}}
 \, \cdot \,
   &\Bigg( \!
     \big( \sqrt{3} + 1 \big) \sqrt{5 - \sqrt{5}}
     - \frac{1}{2} \sqrt{2} \, \big( \sqrt{3} - 1 \big) \big( \sqrt{5} + 1 \big) \, +
   \\
   &\,\,\,\,
     - \, i \, \bigg( \!
     \big( \sqrt{3} - 1 \big) \sqrt{5 - \sqrt{5}}
     + \frac{1}{2} \sqrt{2} \, \big( \sqrt{3} + 1 \big) \big( \sqrt{5} + 1 \big) \!
     \bigg) \!
   \Bigg)
\end{align*}
and:
\begin{align*}
 \cos\paren{\ang{1}}
 = \frac{1}{16} \, \sqrt[3]{\frac{1 + i \sqrt{3}}{2}}
 \, \cdot \,
   &\Bigg( \!
     \big( \sqrt{3} - 1 \big) \sqrt{5 - \sqrt{5}}
     + \frac{1}{2} \sqrt{2} \, \big( \sqrt{3} + 1 \big) \big( \sqrt{5} + 1 \big) \, +
   \\
   &\,\,\,\,
     - \, i \, \bigg( \!
     \big( \sqrt{3} + 1 \big) \sqrt{5 - \sqrt{5}}
     - \frac{1}{2} \sqrt{2} \, \big( \sqrt{3} - 1 \big) \big( \sqrt{5} + 1 \big) \!
     \bigg) \!
   \Bigg) \,\, +
 \\
 + \, \frac{1}{16} \, \sqrt[3]{\frac{1 - i \sqrt{3}}{2}}
 \, \cdot \,
   &\Bigg( \!
     \big( \sqrt{3} - 1 \big) \sqrt{5 - \sqrt{5}}
     + \frac{1}{2} \sqrt{2} \, \big( \sqrt{3} + 1 \big) \big( \sqrt{5} + 1 \big) \, +
   \\
   &\,\,\,\,
     + \, i \, \bigg( \!
     \big( \sqrt{3} + 1 \big) \sqrt{5 - \sqrt{5}}
     - \frac{1}{2} \sqrt{2} \, \big( \sqrt{3} - 1 \big) \big( \sqrt{5} + 1 \big) \!
     \bigg) \!
   \Bigg)
\end{align*}

\subsubsection*{Solution 2}
The complex number $\cos\paren{\ang{1}} + i \sin\paren{\ang{1}}$ is:
\begin{align*}
 \cos\paren{\ang{1}} + i \sin\paren{\ang{1}}
 = e^{-i \frac{\pi}{9}}
 \cdot \big( \! \cos\paren{\ang{21}} &+ i \sin\paren{\ang{21}} \big) = \\
 = \frac{1}{8} \, \sqrt[3]{\frac{1 - i \sqrt{3}}{2}}
 \, \cdot \,
   &\Bigg( \!
     \big( \sqrt{3} - 1 \big) \sqrt{5 - \sqrt{5}}
     + \frac{1}{2} \sqrt{2} \, \big( \sqrt{3} + 1 \big) \big( \sqrt{5} + 1 \big) \, +
   \\
   &\,\,\,\,
     + \, i \, \bigg( \!
     \big( \sqrt{3} + 1 \big) \sqrt{5 - \sqrt{5}}
     - \frac{1}{2} \sqrt{2} \, \big( \sqrt{3} - 1 \big) \big( \sqrt{5} + 1 \big) \!
     \bigg) \!
   \Bigg)
\end{align*}
Then, $\sin\paren{\ang{1}}$ and $\cos\paren{\ang{1}}$
are its imaginary part and its real part, respectively.
